\documentclass[10pt]{extarticle}
\usepackage[a4paper,margin=1.5cm]{geometry}
\usepackage[T1]{fontenc}
\usepackage{amsmath,amssymb,mathtools}
\usepackage{graphicx}
\usepackage[export]{adjustbox}
\usepackage{multicol}
\usepackage{enumitem}
\usepackage{booktabs,array,tabularx}
\setlength{\columnsep}{0.7cm}
\setlength{\parindent}{0pt}
\setlength{\parskip}{3pt}
\pagestyle{empty}
\makeatletter
\newcommand{\qfigdef}[4]{\@namedef{qf@file@#1}{#2}\@namedef{qf@w@#1}{#3}\@namedef{qf@wide@#1}{#4}}
\newcommand{\qf@check}[1]{\@ifundefined{qf@file@#1}{\PackageError{score2}{Unknown figure id #1}{}}{}}
\newcommand{\qf@img}[4][]{\includegraphics[width=#2\dimexpr\@nameuse{qf@w@#4}pt\relax,max width=#3,max height=0.45\textheight,#1]{\@nameuse{qf@file@#4}}}
\newif\ifqf@full
\newif\ifqf@boxed
\newsavebox{\qf@box}
\newcommand{\qf@open}{\global\qf@fulltrue\par\global\qf@boxedtrue
  \setbox\qf@box\vbox\bgroup\hsize\textwidth\linewidth\textwidth}
\newcommand{\qf@emit}[1]{\par\egroup\global\qf@boxedfalse
  \ifdim\dimexpr\ht\qf@box+\dp\qf@box\relax<#1\textheight\noindent\usebox{\qf@box}\par\else\unvbox\qf@box\fi}
\newcommand{\qfullsplit}{\ifqf@boxed\qf@emit{0.85}\fi}
\newcommand{\qf@close}{\ifqf@boxed\qf@emit{0.45}\fi\global\qf@fullfalse}
\newcommand{\qfullbegin}{\end{multicols}\qf@open}
\newcommand{\qfullend}{\qf@close\begin{multicols}{2}}
\newcommand{\qfullstart}{\qf@open}
\newcommand{\qfulldone}{\qf@close}
\newcommand{\qfig}[2][1]{\qf@check{#2}%
  \ifqf@full\qf@blk{#1}{#2}\else
  \ifnum\@nameuse{qf@wide@#2}=1 \expandafter\@firstoftwo\else\expandafter\@secondoftwo\fi
  {\end{multicols}\begin{center}\qf@img{#1}{\textwidth}{#2}\end{center}\begin{multicols}{2}}%
  {\qf@blk{#1}{#2}}\fi}
\newcommand{\qf@blk}[2]{\par\smallskip{\centering\qf@img{#1}{\linewidth}{#2}\par}\smallskip}
\newcommand{\qfigi}[2][1]{\qf@check{#2}\qf@img[valign=c]{#1}{\linewidth}{#2}}
\makeatother
\newcommand{\papersec}[1]{\par\bigskip{\centering\Large\bfseries PAPER-#1\par}\medskip}
\newcommand{\partsec}[2]{\par\medskip\noindent{\large\bfseries PART-#1: #2}\par\smallskip}
\newcommand{\sect}[1]{\par\smallskip\noindent\textbf{#1}\par}
\newcommand{\qn}[1]{\par\smallskip\noindent\textbf{#1.}~}
\newcommand{\paraq}[2]{\par\smallskip\noindent\textbf{Paragraph for Questions #1 and #2}\par}
\newenvironment{qoptions}{\begin{itemize}[leftmargin=1.8em,itemsep=1pt,topsep=2pt,label={}]}{\end{itemize}}
\newcommand{\opt}[2]{\item[(#1)] #2}
\newcolumntype{L}[1]{>{\raggedright\arraybackslash}p{#1}}
\newcommand{\fwstart}{\end{multicols}\par\noindent\begin{minipage}{\textwidth}\small\setlength{\tabcolsep}{4pt}\renewcommand{\arraystretch}{1.5}}
\newcommand{\fwend}{\end{minipage}\par\medskip\begin{multicols}{2}}
% figure registry (generated by score2_tools.py finalize)
\qfigdef{p1_phy_s1_q02_1}{SRG (Enthusiast) - Test 01 (20 Sep 2020)-img/p1_phy_s1_q02_1.png}{105.8}{0}
\qfigdef{p1_phy_s1_q02_2}{SRG (Enthusiast) - Test 01 (20 Sep 2020)-img/p1_phy_s1_q02_2.png}{245.2}{0}
\qfigdef{p1_phy_s1_q03_1}{SRG (Enthusiast) - Test 01 (20 Sep 2020)-img/p1_phy_s1_q03_1.png}{131.8}{0}
\qfigdef{p1_phy_s1_q04_1}{SRG (Enthusiast) - Test 01 (20 Sep 2020)-img/p1_phy_s1_q04_1.png}{119.1}{0}
\qfigdef{p1_phy_s1_q04_2}{SRG (Enthusiast) - Test 01 (20 Sep 2020)-img/p1_phy_s1_q04_2.png}{101.5}{0}
\qfigdef{p1_phy_s2_q06_1}{SRG (Enthusiast) - Test 01 (20 Sep 2020)-img/p1_phy_s2_q06_1.png}{168.2}{0}
\qfigdef{p1_phy_s2_q07_1}{SRG (Enthusiast) - Test 01 (20 Sep 2020)-img/p1_phy_s2_q07_1.png}{122.6}{0}
\qfigdef{p1_phy_s2_q08_1}{SRG (Enthusiast) - Test 01 (20 Sep 2020)-img/p1_phy_s2_q08_1.png}{92.9}{0}
\qfigdef{p1_phy_s2_q09_1}{SRG (Enthusiast) - Test 01 (20 Sep 2020)-img/p1_phy_s2_q09_1.png}{139.8}{0}
\qfigdef{p1_phy_s3_q11_1}{SRG (Enthusiast) - Test 01 (20 Sep 2020)-img/p1_phy_s3_q11_1.png}{131.5}{0}
\qfigdef{p1_phy_s3_q13_1}{SRG (Enthusiast) - Test 01 (20 Sep 2020)-img/p1_phy_s3_q13_1.png}{103.0}{0}
\qfigdef{p1_phy_s3_q14_1}{SRG (Enthusiast) - Test 01 (20 Sep 2020)-img/p1_phy_s3_q14_1.png}{96.9}{0}
\qfigdef{p1_phy_s3_q16_1}{SRG (Enthusiast) - Test 01 (20 Sep 2020)-img/p1_phy_s3_q16_1.png}{78.7}{0}
\qfigdef{p1_phy_s3_q18_1}{SRG (Enthusiast) - Test 01 (20 Sep 2020)-img/p1_phy_s3_q18_1.png}{106.5}{0}
\qfigdef{p1_chem_s1_q04_A}{SRG (Enthusiast) - Test 01 (20 Sep 2020)-img/p1_chem_s1_q04_A.png}{77.3}{0}
\qfigdef{p1_chem_s1_q04_B}{SRG (Enthusiast) - Test 01 (20 Sep 2020)-img/p1_chem_s1_q04_B.png}{77.3}{0}
\qfigdef{p1_chem_s1_q04_C}{SRG (Enthusiast) - Test 01 (20 Sep 2020)-img/p1_chem_s1_q04_C.png}{74.8}{0}
\qfigdef{p1_chem_s1_q04_D}{SRG (Enthusiast) - Test 01 (20 Sep 2020)-img/p1_chem_s1_q04_D.png}{88.2}{0}
\qfigdef{p1_chem_s2_q05_1}{SRG (Enthusiast) - Test 01 (20 Sep 2020)-img/p1_chem_s2_q05_1.png}{245.2}{0}
\qfigdef{p1_chem_s2_q05_2}{SRG (Enthusiast) - Test 01 (20 Sep 2020)-img/p1_chem_s2_q05_2.png}{222.5}{0}
\qfigdef{p1_chem_s3_q11_1}{SRG (Enthusiast) - Test 01 (20 Sep 2020)-img/p1_chem_s3_q11_1.png}{278.0}{1}
\qfigdef{p1_chem_s3_q12_1}{SRG (Enthusiast) - Test 01 (20 Sep 2020)-img/p1_chem_s3_q12_1.png}{299.8}{1}
\qfigdef{p1_chem_s3_q14_1}{SRG (Enthusiast) - Test 01 (20 Sep 2020)-img/p1_chem_s3_q14_1.png}{293.0}{1}
\qfigdef{p1_chem_s3_q17_1}{SRG (Enthusiast) - Test 01 (20 Sep 2020)-img/p1_chem_s3_q17_1.png}{153.4}{0}
\raggedcolumns\makeatletter\@beginparpenalty=10000\makeatother

\begin{document}
\begin{center}{\Large\bfseries SRG (Enthusiast) --- Test 01 (20 Sep 2020)}\end{center}
\begin{multicols}{2}
\partsec{1}{PHYSICS}
\sect{SECTION-I (i)}
\qn{1} A parallel plate capacitor `A' of capacitance 1 $\mu$F is charged to a potential drop 100 volt and then disconnected from the battery. Now the capacitor `A' is completely filled with a dielectric slab of dielectric constant k = 4. Another parallel plate capacitor `B' of capacitance 2 $\mu$F is charged to a potential drop 20 volt and then disconnected from the battery. Now the two capacitors `A' and `B' are connected with each other with opposite polarities. Then the heat dissipated after connecting the capacitors is
\begin{qoptions}
\opt{A}{$1.15\times10^{-3}$ J}
\opt{B}{$1.25\times10^{-3}$ J}
\opt{C}{$1.35\times10^{-3}$ J}
\opt{D}{$1.45\times10^{-3}$ J}
\end{qoptions}
\qn{2} The switch in figure (a) closes when $V_c=\dfrac{2V}{3}$ and opens when $V_c=\dfrac{V}{3}$. The voltmeter reads a voltage as plotted in figure (b). What is the period T of the waveform ? (take $R_1$ = 6R and $R_2$ = 3R; and ln2 = 0.7, if required)
\qfig{p1_phy_s1_q02_1}
\qfig{p1_phy_s1_q02_2}
\begin{qoptions}
\opt{A}{6.3 RC}
\opt{B}{8.4 RC}
\opt{C}{10.5 RC}
\opt{D}{1.4 RC}
\end{qoptions}
\qn{3} The deflection in galvanometer is proportional to current passing through it. In the following circuit, we can record two deflections in galvanometer; one by closing $K_1$ alone and another by closing both $K_1$ \& $K_2$. If one of the deflection is half the other, in one set of reading. If R = 10$\Omega$ and S = 4$\Omega$, then the resistance of the galvanometer is
\qfig{p1_phy_s1_q03_1}
\begin{qoptions}
\opt{A}{$\dfrac{20}{3}\Omega$}
\opt{B}{$\dfrac{40}{3}\Omega$}
\opt{C}{$\dfrac{50}{3}\Omega$}
\opt{D}{$\dfrac{70}{3}\Omega$}
\end{qoptions}
\qn{4} A curved rectangular bar forms a resistor. The curved sides are concentric circular arcs. If $\rho$ is the resistivity of material of bar, $L_0$ is the length of inner arc of radius $r_0$, $r_0+b$ is the radius of outer arc and a is the width of the bar. The electric resistance of the bar across its rectangular ends is
\qfig{p1_phy_s1_q04_1}
\qfig{p1_phy_s1_q04_2}
\begin{qoptions}
\opt{A}{$\dfrac{\rho L_0}{ab}$}
\opt{B}{$\dfrac{\rho L_0}{ar_0\ln\left(1+\dfrac{b}{r_0}\right)}$}
\opt{C}{$\dfrac{2\rho L_0}{ar_0\ln\left(1+\dfrac{b}{r_0}\right)}$}
\opt{D}{$\dfrac{\rho L_0}{2ar_0\ln\left(1+\dfrac{b}{r_0}\right)}$}
\end{qoptions}
\sect{SECTION-I (ii)}
\qn{5} A spherical metallic shell of radius R has a charge Q. This shell is kept in an infinite medium of permittivity $\epsilon$ and resistivity $\rho$. The electrical conduction in the medium follows ohm's law. If the initial charge on the shell is $Q_0$, then
\begin{qoptions}
\opt{A}{Current density as a function of time \& radial distance is given by $\vec J(r,t)=\dfrac{Q_0}{4\pi\epsilon\rho r^2}e^{\frac{-t}{\rho\epsilon}}\cdot\hat r$}
\opt{B}{Current density as a function of time \& radial distance is given by $\vec J(r,t)=\dfrac{Q_0}{2\pi\epsilon\rho r}e^{\frac{-t}{2\rho\epsilon}}\cdot\hat r$}
\opt{C}{Charge on shell as a function of time \& radial distance is given by $Q(r,t)=Q_0e^{\frac{-2t}{\rho\epsilon}}$}
\opt{D}{Charge on shell as a function of time \& radial distance is given by $Q(r,t)=Q_0e^{\frac{-t}{2\rho\epsilon}}$}
\end{qoptions}
\qn{6} For the following circuit, in steady state
\qfig{p1_phy_s2_q06_1}
\begin{qoptions}
\opt{A}{Potential at P is 2V}
\opt{B}{Current through $2\Omega$ is 3A}
\opt{C}{Charge on $C_2$ is 10$\mu$C}
\opt{D}{Charge on $C_3$ is 20$\mu$C}
\end{qoptions}
\qn{7} In the given circuit switch S is closed at t = 0 (you may use $e^{-2}$ = 0.135, $e^{-4}$ = 0.018)
\qfig{p1_phy_s2_q07_1}
\begin{qoptions}
\opt{A}{Time constant is 0.1 ms}
\opt{B}{Charge on the capacitor as a function of time is 400 $(1-e^{-10000t})$ $\mu$C}
\opt{C}{Current through $8\Omega$ resistor at $t=0.4\times10^{-3}$ sec is 0.27A}
\opt{D}{Charge on the capacitor is steady state is 400$\mu$C}
\end{qoptions}
\qn{8} Two conductors AB and CD are made of same conducting material and are connected in parallel as shown. Cross section of AB is more than of CD
\qfig{p1_phy_s2_q08_1}
\begin{qoptions}
\opt{A}{Current density at A, B, C \& D is same}
\opt{B}{Current at A, B, C, D is same}
\opt{C}{Electric field at A, B, C, D is same}
\opt{D}{Drift speed at A, B, C, D is same}
\end{qoptions}
\qn{9} The circuit shown in figure consists of two capacitors of capacitance 2$\mu$F and 4$\mu$F, a resistor of resistance R=$10\Omega$ and an ideal battery of emf 3V. Initially the switch S was open for a long time and the capacitors were fully charged. Then after the switch S is closed,
\qfig{p1_phy_s2_q09_1}
\begin{qoptions}
\opt{A}{Heat produced in the circuit after the switch is closed is 12$\mu$J}
\opt{B}{Energy supplied by the battery after the switch is closed is 24$\mu$J}
\opt{C}{Change in energy stored in the capacitor of capacitance 4$\mu$F after the switch is closed is 16$\mu$J}
\opt{D}{Charge flown through the battery after the switch is closed is 12$\mu$C}
\end{qoptions}
\qn{10} A voltmeter and an ammeter are joined in series to an ideal cell, giving readings V and A respectively. If a resistance equal to resistance of ammeter is now joined parallel to ammeter then:-
\begin{qoptions}
\opt{A}{V will decrease slightly}
\opt{B}{V will increase slightly}
\opt{C}{A will become half of initial value}
\opt{D}{A will become slightly more than half of initial value}
\end{qoptions}
\sect{SECTION-II}
\qn{11} In the following circuit, initially switch S was closed for long time. Switch is opened at t = 0. Find the current (in amperes) through the battery at time RC ln2. Take V = 12 volt, R = 4 ohm, C = 1 micro Farad (take $2^{\frac{-1}{4}}$ = 0.84)
\qfig{p1_phy_s3_q11_1}
\qn{12} The voltage of clouds is $4\times10^6$ V with respect to ground. In a lightning strike last 100 ms, a charge of 4 C is delivered to the ground. The average power of lightning strike is (assume complete discharges) (in MW)
\qn{13} Two circular rings of identical radii and resistance $36\Omega$ are placed as shown in the figure. Conducting joints are made at points A and B and cell of emf 20 V is connected between these two points. What is the power (in W) delivered by the cell? {[}$C_1$ and $C_2$ are the centres of the two rings{]}
\qfig{p1_phy_s3_q13_1}
\qn{14} The figure below shows a wheatstone bridge with resistors P and Q having almost equal resistance. When R = $400\Omega$. The bridge is in balanced condition. If on interchanging P and Q, the bridge is again balanced for R = $405\Omega$, then the value of X is
\qfig{p1_phy_s3_q14_1}
\qn{15} Consider a cylindrical conductor of length L and area of cross-section A. The specific conductivity varies as $\sigma(x)=\sigma_0\dfrac{L}{\sqrt{x}}$ where x is the distance along the axis of the cylinder from one of its ends. The resistance of the system along the cylindrical axis is $\dfrac{x\sqrt{L}}{6A\sigma_0}$. Find x
\qn{16} The V-I graphs for a conductor at temperature $T_1$ \& $T_2$ are shown. Let $R_0$ be the resistance at $0^\circ$C, $T_1$ \& $T_2$ be temperature in $^\circ$C and $\alpha$ is the thermal coefficient of resistance in /$^\circ$C. If $T_2-T_1=\dfrac{1}{n\alpha}$ find n. (Given $R_0=\sqrt3\,\Omega$)
\qfig{p1_phy_s3_q16_1}
\qn{17} Through a coil of resistance R charge is flowing in such a manner that in every time interval $t_0$ the amount of charge flow is reduced by a factor of 1/2. If a total of $Q_0$ charge flows through the coil in very long time, then the total heat generated in the coil will be $\dfrac{Q_0^2R\ln2}{nt_0}$. Find n
\qn{18} Distance between consecutive metallic plates is d and each plate has an area A. The capacitance between terminals A \& B is $\dfrac{n\epsilon_0A}{d}$. Find n
\qfig{p1_phy_s3_q18_1}
\partsec{2}{CHEMISTRY}
\sect{SECTION-I (i)}
\qn{1} Which of the following is not a semisynthetic polymer?
\begin{qoptions}
\opt{A}{Cis-polyisoprene}
\opt{B}{Cellulose nitrate}
\opt{C}{Cellulose acetate}
\opt{D}{Vulcanised rubber}
\end{qoptions}
\qn{2} The percentage of sulphur in the organic compound, when 0.2595 g of a sulphur containing organic compound in a quantitative analysis by carius method yielded 0.35 g of barium sulphate is
\begin{qoptions}
\opt{A}{14.52\%}
\opt{B}{16.52\%}
\opt{C}{18.52\%}
\opt{D}{19.52\%}
\end{qoptions}
\qn{3} GUTTA PERCHA is
\begin{qoptions}
\opt{A}{Cis-polyisoprene}
\opt{B}{Trans polyisoprene}
\opt{C}{Chloroprene}
\opt{D}{Diphenyl guanidine}
\end{qoptions}
\qn{4} The structural formula of the compound (A) having molecular formula C$_6$H$_{11}$N, which is resolvable, dissolves in dilute aqueous HCl and also produces N$_2$ on reaction with nitrous acid is
\begin{qoptions}
\opt{A}{\qfigi{p1_chem_s1_q04_A}}
\opt{B}{\qfigi{p1_chem_s1_q04_B}}
\opt{C}{\qfigi{p1_chem_s1_q04_C}}
\opt{D}{\qfigi{p1_chem_s1_q04_D}}
\end{qoptions}
\sect{SECTION-I (ii)}
\qn{5} The correct statement(s) about the following sugars X and Y is(are)
\qfig{p1_chem_s2_q05_1}
{\centering (X)\par}
\qfig{p1_chem_s2_q05_2}
{\centering (Y)\par}
\begin{qoptions}
\opt{A}{X is a reducing sugar and Y is a non-reducing sugar}
\opt{B}{X is a non-reducing sugar and Y is a reducing sugar}
\opt{C}{The glucosidic linkages in X and Y are $\alpha$ and $\beta$ respectively}
\opt{D}{The glucosidic linkages in X and Y are $\beta$ and $\alpha$ respectively}
\end{qoptions}
\qn{6} Find the correct statement(s)
\begin{qoptions}
\opt{A}{Methyl amine in water reacts with FeCl$_3$ to precipitate hydrated ferric oxide}
\opt{B}{Diazonium salt of aromatic amines are more stable than those of aliphatic amines}
\opt{C}{Aniline does not undergo friedal-craft's reaction}
\opt{D}{Gabriel pthalimide synthesis is preffered for 1$^\circ$ amine synthesis}
\end{qoptions}
\qn{7} Nylon is used as a general name for all synthetic fibres forming polyamides. Identify the correct statement about Nylon.
\begin{qoptions}
\opt{A}{Nylon-610 is a condensation polymerisation of hexamethylenediamine and sebacic acid}
\opt{B}{Natural rubber is cis polyisoprene}
\opt{C}{Nylon-6 is produced by prolonged heating of caprolactum}
\opt{D}{PHBV is a biodegradable polymer}
\end{qoptions}
\qn{8} Which of the following is/are correct about aniline ?
\begin{qoptions}
\opt{A}{Lassaigne's extract of its solution on treatment with FeSO$_4$ / H$_2$SO$_4$ give blue colouration}
\opt{B}{It gives Fehling's test}
\opt{C}{It gives carbylamine test}
\opt{D}{31 g of aniline on combustion gives 88 g of CO$_2$}
\end{qoptions}
\qn{9} Which of the following reagents is/are suitable to distinguish between the given amines ?\\
CH$_3$CH$_2$NH$_2$, CH$_3$ $-$ NH $-$ CH$_2$CH$_3$ and (CH$_3$)$_3$N
\begin{qoptions}
\opt{A}{NaNO$_2$ / HCl}
\opt{B}{CS$_2$, HgCl$_2$, $\Delta$}
\opt{C}{CHCl$_3$ / KOH}
\opt{D}{(CO$_2$C$_2$H$_5$)$_2$}
\end{qoptions}
\qn{10} Which of the following polymers are thermoplastic ?
\begin{qoptions}
\opt{A}{Teflon}
\opt{B}{Natural rubber}
\opt{C}{Neoprene}
\opt{D}{Polysteyrene}
\end{qoptions}
\qfullbegin
\sect{SECTION-II}
\qn{11} What is the net charge of the below given compound at pH = 8 ?
\qfig{p1_chem_s3_q11_1}
\qn{12} X = Number of compounds give 3$^\circ$ alcohol when reacts with HNO$_2$ ?
\qfig{p1_chem_s3_q12_1}
\qfullend
\qn{13} The number of tripeptides synthesized using amino acids: glycine, alanine and valine are \underline{\hspace{1cm}}
\qfullbegin
\qn{14} X = Number of compounds give positive isocyanide test.
\qfig{p1_chem_s3_q14_1}
\qfullend
\qn{15} What is PI of the following compound\\
\begin{tabular}{@{}l@{}}
\hspace{1.2em}O\\
\hspace{1.2em}$\|$\\
HO $-$ C $-$ CH$_2$ $-$ CH $-$ COOH\\
\hspace{0.2em}(pK$_\mathrm{a}$=4) \hspace{2.2em}$|$ \hspace{1.3em}(pK$_\mathrm{a}$=2)\\
\hspace{6.3em}$+$NH$_3$\\
\hspace{6.0em}(pK$_\mathrm{a}$=9)
\end{tabular}
\qn{16} RCONH$_2$ + Br$_2$ + KOH $\longrightarrow$ RNH$_2$ + KBr + K$_2$CO$_3$ + H$_2$O\\
Total moles of KOH consumed per mole RCONH$_2$ in this reaction will be \underline{\hspace{1cm}}
\qn{17} Molecular weight of (A) increase by \underline{\hspace{1cm}}
\qfig{p1_chem_s3_q17_1}
\qn{18} How many of the following polymers are addition polymers ?\\
(i) Polyvinyl chloride \hfill (ii) Terylene\\
(iii) Teflon \hfill (iv) Neoprene\\
(v) Buna-S \hfill (vi) Nylon-6,6\\
(vii) Natural rubber \hfill (viii) Bakelite
\partsec{3}{MATHEMATICS}
\sect{SECTION-I (i)}
\qn{1} Let $U_{10}=\displaystyle\int_0^{\pi/2}x\cdot\sin^{10}x\,dx$, then the value of $\left(\dfrac{100U_{10}-1}{U_8}\right)$ is
\begin{qoptions}
\opt{A}{90}
\opt{B}{100}
\opt{C}{110}
\opt{D}{80}
\end{qoptions}
\qn{2} Let $f(x)$ be an even function such that $\displaystyle\int_0^\infty f(x)\,dx=\dfrac{\pi}{2}$, then the value of definite integral $\displaystyle\int_0^\infty f\left(x-\dfrac1x\right)dx$ is
\begin{qoptions}
\opt{A}{$\pi$}
\opt{B}{$\dfrac{\pi}{2}$}
\opt{C}{$\dfrac{3\pi}{2}$}
\opt{D}{$2\pi$}
\end{qoptions}
\qn{3} The integral $\displaystyle\int\dfrac{7x^4+5x^2+4x}{\sqrt{x^3+x+1}}\,dx$ is
\begin{qoptions}
\opt{A}{$\left(x^7+x^5+x^4\right)^{1/2}+c$}
\opt{B}{$2\left(x^7+x^5+x^4\right)^{1/2}+c$}
\opt{C}{$\left(x^3+x+1\right)^{-1/2}+c$}
\opt{D}{$2\left(x^3+x+1\right)^{-1/2}+c$}
\end{qoptions}
\qn{4} If $\displaystyle\int_{-1}^{1}e^{\left(x^5-5x^4+10x^3-10x^2+5x+10+4\ln|x-1|\right)}dx=\dfrac{e^p-1}{5e^q}$ ; then p + q equals:
\begin{qoptions}
\opt{A}{23}
\opt{B}{32}
\opt{C}{45}
\opt{D}{55}
\end{qoptions}
\sect{SECTION-I (ii)}
\qn{5} The function `f ' is continuous and has the property $f\{f(x)\}=1-x\ \forall\ x\in[0,1]$ and $J=\displaystyle\int_0^1f(x)\,dx$, then
\begin{qoptions}
\opt{A}{$f\left(\dfrac14\right)+f\left(\dfrac34\right)=1$}
\opt{B}{$\displaystyle\int_0^1\dfrac{48x^6-32x^3}{\left(3x^3+1\right)^3}dx=J$}
\opt{C}{$\dfrac12\displaystyle\int_0^\pi\left|\dfrac{1}{\sqrt3}\sin x+\dfrac{1}{\sqrt6}\cos x\right|dx=J$}
\opt{D}{$\displaystyle\int_0^\pi\dfrac{\sin x\,dx}{\left(\sin x+\cos x\right)^3}=J$}
\end{qoptions}
\qn{6} If $I_{m,n}=\displaystyle\int\cos^mx\sin nx\,dx$, then $7I_{4,3}-4I_{3,2}$ is equal to
\begin{qoptions}
\opt{A}{Constant}
\opt{B}{$-\cos^2x+c$}
\opt{C}{$-\cos^4x\cos3x+c$}
\opt{D}{$\cos7x-\cos4x+c$}
\end{qoptions}
\qn{7} If
{\small\[\begin{gathered}\int\dfrac{x^{-7/6}-x^{5/6}}{x^{1/3}\cdot\left(x^2+x+1\right)^{1/2}-x^{1/2}\left(x^2+x+1\right)^{1/3}}\,dx\\=-A\left\{\dfrac{z^3}{3}+\dfrac{3z^B}{2}+\dfrac{3z^C}{D}+\log|z|\right\}+k\end{gathered}\]}
where $z=\left\{x+\dfrac1x+1\right\}^{\frac16}-1$, then
\begin{qoptions}
\opt{A}{A = 6}
\opt{B}{B + C = 4}
\opt{C}{B + C = 3}
\opt{D}{B + C + D = 5}
\end{qoptions}
\qn{8} If $\displaystyle\int\dfrac{\sqrt{4+x^2}}{x^6}\,dx=\dfrac{\alpha\left(4+x^2\right)^{3/2}\left(\beta x^2-6\right)}{x^5}+C$ , then
\begin{qoptions}
\opt{A}{$\alpha=\dfrac{1}{120}$}
\opt{B}{$\beta=1$}
\opt{C}{$\alpha=\dfrac{-1}{120}$}
\opt{D}{$\beta=-1$}
\end{qoptions}
\qn{9} If $\displaystyle\int\dfrac{dx}{x^2+ax+1}=f\left(g(x)\right)+c$ , then
\begin{qoptions}
\opt{A}{$f(x)$ is inverse trigonometric function for $|a|<2$}
\opt{B}{$f(x)$ is logarithmic function for $|a|<2$}
\opt{C}{$f(x)$ is quadratic function for $|a|>2$}
\opt{D}{$g(x)$ is rational function for $|a|<2$}
\end{qoptions}
\qn{10} If $I=\displaystyle\int\dfrac{\left(1+\log_ex\right)^2dx}{1+\log_ex^{x+1}+\left(\log_ex^{\sqrt x}\right)^2}=f\left(g(x)\right)+c$. If $f(x)=\ln x$, then
\begin{qoptions}
\opt{A}{$g(1)=1$}
\opt{B}{$f(1)=1$}
\opt{C}{Domain of $g(x)$ is $x\in R^+$}
\opt{D}{$g(x)$ is periodic with period 1}
\end{qoptions}
\sect{SECTION-II}
\qn{11} If $\displaystyle\int\dfrac{e^{\sin\theta}}{\cos^2\theta}\left\{\theta\cos^3\theta-\sin\theta\right\}d\theta=f(\theta)\left[\theta-\sec\theta\right]+K$, where `K' is constant of integration, then largest integral value of $f(\theta)$ is \underline{\hspace{1cm}}
\qn{12} {\small\[\begin{gathered}\int\dfrac{2(x+16)^{1/2}+(x+16)^{3/4}}{x(x+16)^{5/4}}dx\\=\dfrac1p\left[\ln\left|(x+16)^{1/4}-2\right|-\dfrac12\ln\left((x+16)^{1/2}+4\right)\right.\\\left.-\tan^{-1}\dfrac{(x+16)^{1/4}}{2}\right]+c.\end{gathered}\]}
Where c is constant of integration then $p^2$ is -
\qn{13} If $\displaystyle\int\sqrt{\dfrac{x^{2009}}{2x^{2008}-x^{2009}}}\,dx=\dfrac p2\sin^{-1}\sqrt{\dfrac x2}+\dfrac q2\sqrt{2x-x^2}+c$ , where c is constant of integration then $|p-q|$ is
\qn{14} If $\displaystyle\int\tan^2x\,\sin^{-1}\left(\tan x-x\right)dx=G(x)+c$, where c is constant of integration then G(0) is
\qn{15} $\displaystyle\int\left(\dfrac{\sec^2x}{\tan^{-1}x}-\dfrac{\tan x}{\left(\tan^{-1}x\right)^2\left(1+x^2\right)}\right)dx=\dfrac{\tan x}{f(x)}+c$ , where c is constant of integration, then $1/f\,'(2)$ is....
\qn{16} If $f(\pi)=2$ and $\displaystyle\int_0^\pi\left(f(x)+f''(x)\right)\sin x\,dx=5$ , then f(0) is
\qn{17} If $2f(x)+f(-x)=\dfrac1x\sin\left(x-\dfrac1x\right)$ then $\displaystyle\int_{1/e}^{e}f(x)\,dx=$ \underline{\hspace{1cm}}
\qn{18} If $f(x)=\begin{vmatrix}\cos x&e^{x^2}&2x\cos^2\dfrac x2\\x^2&\sec x&\sin x+x^3\\1&2&x+\tan x\end{vmatrix}$ then the value of $\displaystyle\int_{-\pi/2}^{\pi/2}\left(x^2+1\right)\left\{f(x)+f''(x)\right\}dx$ is \underline{\hspace{1cm}}
\end{multicols}
\end{document}
